% GENERATED FILE. Edit canonical lessons, then run npm run generate:books.
% canonical-source-hash: fnv1a64:4eca8e148613788b
% canonical-lessons: UR-C494-intarvyu, UR-C494-rozgar, UR-C494-asami, UR-C494-bartaraf, UR-C494-bharti-karna

\chapter{Looking for Work}
\label{ch:ur-looking-for-work}
\hlchaptermodality{494}

\begin{chapteropening}
\textbf{By the end of this chapter:} \emph{I can talk about looking for a job and getting one.}

Complete the last lesson of chapter 494: I can talk about looking for a job and getting one.
\end{chapteropening}

\section[\texorpdfstring{\ur{انٹرویو} (inṭarvyū)}{انٹرویو (inṭarvyū)}]{\ur{انٹرویو} (inṭarvyū) — an interview (for a job)}

\label{lesson:UR-C494-intarvyu}



\begin{quote}
Before the new one: say the Urdu for a fruit seller, then the Urdu for a greengrocer.
\end{quote}

\subsection*{You'll want to know: \ur{انٹرویو}}
\textbf{\ur{انٹرویو}} — \emph{inṭarvyū} — \textquotedblleft{}an interview (for a job)\textquotedblright{}.

\begin{quote}
\ur{کل} \ur{میرا} \ur{ایک} \ur{کمپنی} \ur{میں} \ur{انٹرویو} \ur{ہے۔}
\end{quote}

\begin{quote}
\emph{kal merā ek kampanī meṅ inṭarvyū hai.}
\end{quote}

\begin{quote}
Tomorrow I have an interview at a company.
\end{quote}

Masculine.

\begin{cousinweb}
\textbf{\ur{انٹرویو}} (\emph{inṭarvyū}) is borrowed from English \emph{interview}.
\end{cousinweb}

\begin{grammarlens}[title={What you've built}]
One more word for describing what you lived and felt.
\end{grammarlens}

\subsection*{Your turn}
\begin{itemize}
\raggedright
  \item \emph{Say it:} \emph{inṭarvyū}
  \item \emph{Say it:} \emph{inṭarvyū}, once more
  \item \emph{Say it:} \emph{kal merā ek kampanī meṅ inṭarvyū hai.} — Tomorrow I have an interview at a company.
  \item \emph{Recall:} say the Urdu for a fruit seller, then the Urdu for a greengrocer, then say \emph{inṭarvyū} again
\end{itemize}

\begin{tcolorbox}[breakable,colback=teal!4,colframe=teal!35!black,title={Before you move on}]
What does \ur{انٹرویو} mean? (\textquotedblleft{}An interview (for a job)\textquotedblright{}.) Say it once more.
\end{tcolorbox}
\section[\texorpdfstring{\ur{روزگار} (rozgār)}{روزگار (rozgār)}]{\ur{روزگار} (rozgār) — employment; a livelihood}

\label{lesson:UR-C494-rozgar}



\begin{quote}
Before the new one: say the Urdu for a greengrocer, then the Urdu for an interview (for a job).
\end{quote}

\subsection*{You'll want to know: \ur{روزگار}}
\textbf{\ur{روزگار}} — \emph{rozgār} — \textquotedblleft{}employment; a livelihood\textquotedblright{}.

\begin{quote}
\ur{اس} \ur{شہر} \ur{میں} \ur{روزگار} \ur{کم} \ur{ہے۔}
\end{quote}

\begin{quote}
\emph{is shahr meṅ rozgār kam hai.}
\end{quote}

\begin{quote}
There is little work to be had in this city.
\end{quote}

Masculine.

\begin{cousinweb}
\textbf{\ur{روزگار}} (\emph{rozgār}) is usually traced to Persian \emph{rōzgār}, \textquotedblleft{}the days, the times\textquotedblright{}, built on \emph{rōz}, \textquotedblleft{}a day\textquotedblright{}; the sense \textquotedblleft{}a livelihood\textquotedblright{} is the one Urdu uses.
\end{cousinweb}

\begin{grammarlens}[title={What you've built}]
One more word for describing what you lived and felt.
\end{grammarlens}

\subsection*{Your turn}
\begin{itemize}
\raggedright
  \item \emph{Say it:} \emph{rozgār}
  \item \emph{Say it:} \emph{rozgār}, once more
  \item \emph{Say it:} \emph{is shahr meṅ rozgār kam hai.} — There is little work to be had in this city.
  \item \emph{Recall:} say the Urdu for a greengrocer, then the Urdu for an interview (for a job), then say \emph{rozgār} again
\end{itemize}

\begin{tcolorbox}[breakable,colback=teal!4,colframe=teal!35!black,title={Before you move on}]
What does \ur{روزگار} mean? (\textquotedblleft{}Employment; a livelihood\textquotedblright{}.) Say it once more.
\end{tcolorbox}
\section[\texorpdfstring{\ur{اسامی} (asāmī)}{اسامی (asāmī)}]{\ur{اسامی} (asāmī) — a vacancy (a post to fill)}

\label{lesson:UR-C494-asami}



\begin{quote}
Before the new one: say the Urdu for an interview (for a job), then the Urdu for employment.
\end{quote}

\subsection*{You'll want to know: \ur{اسامی}}
\textbf{\ur{اسامی}} — \emph{asāmī} — \textquotedblleft{}a vacancy (a post to fill)\textquotedblright{}.

\begin{quote}
\ur{دفتر} \ur{میں} \ur{ایک} \ur{اسامی} \ur{خالی} \ur{ہے۔}
\end{quote}

\begin{quote}
\emph{daftar meṅ ek asāmī khālī hai.}
\end{quote}

\begin{quote}
There is a vacancy in the office.
\end{quote}

Feminine.

\begin{cousinweb}
\textbf{\ur{اسامی}} (\emph{asāmī}) is usually traced to Arabic \emph{asāmī}, \textquotedblleft{}names\textquotedblright{}, the plural of \emph{ism}, \textquotedblleft{}a name\textquotedblright{}; the sense \textquotedblleft{}a post to be filled\textquotedblright{} came later, and the sources we use do not show how it arose.
\end{cousinweb}

\begin{grammarlens}[title={What you've built}]
One more word for describing what you lived and felt.
\end{grammarlens}

\subsection*{Your turn}
\begin{itemize}
\raggedright
  \item \emph{Say it:} \emph{asāmī}
  \item \emph{Say it:} \emph{asāmī}, once more
  \item \emph{Say it:} \emph{daftar meṅ ek asāmī khālī hai.} — There is a vacancy in the office.
  \item \emph{Recall:} say the Urdu for an interview (for a job), then the Urdu for employment, then say \emph{asāmī} again
\end{itemize}

\begin{tcolorbox}[breakable,colback=teal!4,colframe=teal!35!black,title={Before you move on}]
What does \ur{اسامی} mean? (\textquotedblleft{}A vacancy (a post to fill)\textquotedblright{}.) Say it once more.
\end{tcolorbox}
\section[\texorpdfstring{\ur{برطرف} (bartaraf)}{برطرف (bartaraf)}]{\ur{برطرف} (bartaraf) — dismissed from a job}

\label{lesson:UR-C494-bartaraf}



\begin{quote}
Before the new one: say the Urdu for employment, then the Urdu for a vacancy.
\end{quote}

\subsection*{You'll want to know: \ur{برطرف}}
\textbf{\ur{برطرف}} — \emph{bartaraf} — \textquotedblleft{}dismissed from a job\textquotedblright{}.

\begin{quote}
\ur{کمپنی} \ur{نے} \ur{دو} \ur{ملازموں} \ur{کو} \ur{برطرف} \ur{کیا۔}
\end{quote}

\begin{quote}
\emph{kampanī ne do mulāzimoṅ ko bartaraf kiyā.}
\end{quote}

\begin{quote}
The company dismissed two employees.
\end{quote}

\begin{cousinweb}
\textbf{\ur{برطرف}} (\emph{bartaraf}) is usually traced to Persian \emph{bar-taraf}, \textquotedblleft{}set aside\textquotedblright{}, built from \emph{bar}, \textquotedblleft{}on, to\textquotedblright{}, and \emph{taraf}, \textquotedblleft{}a side\textquotedblright{}, the word Urdu also has as \textbf{\ur{طرف}} (\emph{taraf}); Urdu uses it for someone put out of a job.
\end{cousinweb}

\begin{grammarlens}[title={What you've built}]
One more word for describing what you lived and felt.
\end{grammarlens}

\subsection*{Your turn}
\begin{itemize}
\raggedright
  \item \emph{Say it:} \emph{bartaraf}
  \item \emph{Say it:} \emph{bartaraf}, once more
  \item \emph{Say it:} \emph{kampanī ne do mulāzimoṅ ko bartaraf kiyā.} — The company dismissed two employees.
  \item \emph{Recall:} say the Urdu for employment, then the Urdu for a vacancy, then say \emph{bartaraf} again
\end{itemize}

\begin{tcolorbox}[breakable,colback=teal!4,colframe=teal!35!black,title={Before you move on}]
What does \ur{برطرف} mean? (\textquotedblleft{}Dismissed from a job\textquotedblright{}.) Say it once more.
\end{tcolorbox}
\section[\texorpdfstring{\ur{بھرتی} \ur{کرنا} (bhartī karnā)}{بھرتی کرنا (bhartī karnā)}]{\ur{بھرتی} \ur{کرنا} (bhartī karnā) — to recruit; to take on (staff)}

\label{lesson:UR-C494-bharti-karna}



\begin{quote}
Before the new one: say the Urdu for a vacancy, then the Urdu for dismissed from a job.
\end{quote}

\subsection*{You'll want to know: \ur{بھرتی} \ur{کرنا}}
\textbf{\ur{بھرتی} \ur{کرنا}} — \emph{bhartī karnā} — \textquotedblleft{}to recruit; to take on (staff)\textquotedblright{}.

\begin{quote}
\ur{کارخانے} \ur{نے} \ur{دس} \ur{نئے} \ur{مزدور} \ur{بھرتی} \ur{کیے۔}
\end{quote}

\begin{quote}
\emph{kārkhāne ne das naye mazdūr bhartī kiye.}
\end{quote}

\begin{quote}
The factory took on ten new labourers.
\end{quote}

Past, with \ur{نے} (he or she …): \ur{اس} \ur{نے} \ur{دو} \ur{آدمی} \ur{بھرتی} \ur{کیے}, \emph{us ne do ādmī bhartī kiye}.

\begin{cousinweb}
\textbf{\ur{بھرتی} \ur{کرنا}} (\emph{bhartī karnā}) is built from \emph{bhartī}, \textquotedblleft{}filling up; an intake\textquotedblright{}, on the stem of \emph{bharnā}, \textquotedblleft{}to fill\textquotedblright{}, which CDIAL lists under Sanskrit \emph{bharati}, \textquotedblleft{}bears, fills\textquotedblright{} (CDIAL 9397), and \emph{karnā}, \textquotedblleft{}to do\textquotedblright{}.
\end{cousinweb}

\begin{grammarlens}[title={What you've built}]
One more word for describing what you lived and felt.
\end{grammarlens}

\subsection*{Your turn}
\begin{itemize}
\raggedright
  \item \emph{Say it:} \emph{bhartī karnā}
  \item \emph{Say it:} \emph{bhartī karnā}, once more
  \item \emph{Say it:} \emph{kārkhāne ne das naye mazdūr bhartī kiye.} — The factory took on ten new labourers.
  \item \emph{Recall:} say the Urdu for a vacancy, then the Urdu for dismissed from a job, then say \emph{bhartī karnā} again
\end{itemize}

\begin{tcolorbox}[breakable,colback=teal!4,colframe=teal!35!black,title={Before you move on}]
What does \ur{بھرتی} \ur{کرنا} mean? (\textquotedblleft{}To recruit; to take on (staff)\textquotedblright{}.) Say it once more.
\end{tcolorbox}
